% Kapitel 7 — Budget: 14 Seiten — Schreiben: 30.11.–11.12. (fertig 11.12., interner Meilenstein)
% Vormessungen und Studie (Beschluss 26.10.); Tabellen statt Prosa, je Person eine Zeile.

\chapter{Ergebnisse}
\label{chap:ergebnisse}

\section{Ergebnisse der Vormessungen}% 2 S.
% Inhalt: aus Kap. 5 hierher gezogen (Beschluss 26.10.); die fünf wichtigsten
% Messungen ausführlich, die übrigen elf nur kurz (warum, wie, Ergebnis) in einem
% gemeinsamen Abschnitt; volle Tabellen in den Anhang.
\label{sec:erg-vormessungen}

\section{Durchführung und Teilnehmende}% 1,5 S.
% Inhalt: was bei der Durchführung vom Plan abwich (Abbrüche, hängende Läufe,
% fehlende Aufnahmen, Dauer).
\label{sec:durchfuehrung}

\section{Verständniszuwachs unter D gegen C}% 3 S.
% Inhalt: mit Interviewzitaten.
\label{sec:verstaendniszuwachs}

\section{Erklärungen}% 2 S.
\label{sec:erklaerungen}

\section{Eingriffe in den Kontext und Kontextqualität}% 3 S.
% Inhalt: Frage 1 des Studienplans; Nutzung der Oberfläche, SUS.
\label{sec:eingriffe}

\section{Ergebnis des Agenten gegen die Prüfliste, C gegen D}% 2,5 S.
% Inhalt: Scheiternskriterium Exposé 4.4.
\label{sec:agentenergebnis}
